%-------------------------
% Özgeçmiş (LaTeX) - Türkçe Versiyon
% Yazar : Yiğit Arda Kıdıman
%-------------------------

\documentclass[letterpaper,10.5pt]{article}

\usepackage{latexsym}
\usepackage[empty]{fullpage}
\usepackage{titlesec}
\usepackage{marvosym}
\usepackage[usenames,dvipsnames]{color}
\usepackage{verbatim}
\usepackage{enumitem}
\usepackage[hidelinks]{hyperref}
\usepackage{fancyhdr}
\usepackage[turkish]{babel}
\usepackage{tabularx}
\input{glyphtounicode}

%----------FONT SEÇENEKLERİ----------
\usepackage[sfdefault]{Inter}
\usepackage[T1]{fontenc}

\pagestyle{fancy}
\fancyhf{}
\fancyfoot{}
\renewcommand{\headrulewidth}{0pt}
\renewcommand{\footrulewidth}{0pt}

% Kenar Boşlukları
\addtolength{\oddsidemargin}{-0.5in}
\addtolength{\evensidemargin}{-0.5in}
\addtolength{\textwidth}{1in}
\addtolength{\topmargin}{-.5in}
\addtolength{\textheight}{1.0in}

\urlstyle{same}

\raggedbottom
\raggedright
\setlength{\tabcolsep}{0in}

% Bölüm Formatı
\titleformat{\section}{
  \vspace{-5pt}\scshape\raggedright\large\bfseries
}{}{0em}{}[\color{black}\titlerule \vspace{-4pt}]

\pdfgentounicode=1

%-------------------------
% Özel Komutlar
\newcommand{\resumeItem}[1]{
  \item\small{
    {#1 \vspace{-2pt}}
  }
}

\newcommand{\resumeSubheading}[4]{
  \vspace{-2pt}\item
    \begin{tabular*}{0.97\textwidth}[t]{l@{\extracolsep{\fill}}r}
      \textbf{#1} & #2 \\
      \textit{\small#3} & \textit{\small #4} \\
    \end{tabular*}\vspace{-6pt}
}

\newcommand{\resumeProjectHeading}[2]{
    \item
    \begin{tabular*}{0.97\textwidth}{l@{\extracolsep{\fill}}r}
      \small#1 & #2 \\
    \end{tabular*}\vspace{-6pt}
}

\newcommand{\resumeSubHeadingListStart}{\begin{itemize}[leftmargin=0.15in, label={}]}
\newcommand{\resumeSubHeadingListEnd}{\end{itemize}}
\newcommand{\resumeItemListStart}{\begin{itemize}}
\newcommand{\resumeItemListEnd}{\end{itemize}\vspace{-5pt}}

%-------------------------------------------
%%%%%%  ÖZGEÇMİŞ BURADA BAŞLAR  %%%%%%%%%%%%%%%%%%%%%%%%%%%%

\begin{document}

%----------BAŞLIK----------
\begin{center}
    \textbf{\Huge \scshape Yiğit Arda Kıdıman} \\ \vspace{2pt}
    \small Yazılım Mühendisliği Öğrencisi \& Web Geliştirici \\ \vspace{2pt}
    \href{mailto:yigitardakidiman@gmail.com}{\underline{yigitardakidiman@gmail.com}} $|$ 
    \href{https://kidiman.com}{\underline{kidiman.com}} $|$ 
    \href{https://linkedin.com/in/yigitardakidiman}{\underline{linkedin.com/in/yigitardakidiman}} $|$ 
    \href{https://github.com/yigitardakidiman}{\underline{github.com/yigitardakidiman}}
\end{center}

%-----------EĞİTİM-----------
\section{Eğitim}
  \resumeSubHeadingListStart
    \resumeSubheading
      {OSTİM Teknik Üniversitesi}{09/2024 -- Günümüz}
      {Yazılım Mühendisliği (Lisans) $|$ Genel Not Ortalaması (GPA): 3.41 / 4.00}{Ankara, Türkiye}

  \resumeSubHeadingListEnd

%-----------TEKNİK YETENEKLER-----------
\section{Teknik Yetenekler}
 \begin{itemize}[leftmargin=0.15in, label={}]
    \small{\item{
     \textbf{Programlama Dilleri}{: Python, JavaScript, TypeScript, C, HTML5, CSS3} \\
     \textbf{Frontend}{: React.js, Next.js, Three.js, Tailwind CSS} \\
     \textbf{Backend}{: Django, Django REST Framework, REST API'ler} \\
     \textbf{Araçlar \& Veritabanları}{: Supabase, Git, GitHub, Vite, VS Code}
    }}
 \end{itemize}

%-----------DENEYİM VE LİDERLİK-----------
\section{Deneyim \& Liderlik}
  \resumeSubHeadingListStart

    \resumeSubheading
      {Volinor Savunma ve Teknoloji A.Ş.}{06/2026 -- Günümüz}
      {Yazılım Geliştirici (Yarı Zamanlı: 07/2026 -- Günümüz $|$ Stajyer: 06/2026 -- 07/2026)}{Ankara, Türkiye}
      \resumeItemListStart
        \resumeItem{3D model veri yönetimi, ürün katalog listelemeleri, kullanıcı kimlik doğrulama ve e-posta servisleri için \textbf{Django REST API uç noktaları} geliştirdim ve bakımını yaptım.}
        \resumeItem{Şirket ürünlerini web üzerinden sergilemek amacıyla \textbf{React ve Three.js} kullanarak etkileşimli 3D model görüntüleme arayüzleri geliştirdim.}
        \resumeItem{Versiyon kontrolü, kod inceleme süreçleri ve hata çözümlerinde geliştirme ekibiyle Git üzerinden iş birliği yaptım.}
      \resumeItemListEnd

    \resumeSubheading
      {Huawei Student Developers (HSD) OSTIMTech}{10/2025 -- Günümüz}
      {Proje Komite Başkanı (10/2026 -- Günümüz) $|$ Çekirdek Ekip Üyesi (10/2025 -- 10/2026)}{Ankara, Türkiye}
      \resumeItemListStart
        \resumeItem{\textbf{Proje Komite Başkanı} olarak öğrenci yazılım inisiyatiflerini koordine ediyor, topluluk içi teknik etkinlik ve atölye çalışmalarını organize ediyorum.}
        \resumeItem{28+ çekirdek ekip üyesinin kulüp içi görev dağılımını ve takibini kolaylaştırmak için \textbf{React ve Supabase} ile kulüp yönetim paneli geliştirdim.}
        \resumeItem{\textbf{60+ katılımcılı} OSTİMTech üniversite hackathon'unun organizasyon süreçlerine ve teknik koordinasyonuna destek sağladım.}
      \resumeItemListEnd

  \resumeSubHeadingListEnd

%-----------ÖNE ÇIKAN PROJELER-----------
\section{Öne Çıkan Projeler}
    \resumeSubHeadingListStart
      \resumeProjectHeading
          {\textbf{GitFetch -- README Oluşturucu} $|$ \emph{React, TypeScript, Tailwind CSS}}{\href{https://gitfetch.kidiman.com}{Canlı Demo} $|$ \href{https://github.com/yigitardakidiman/gitfetch-readme-generator}{GitHub}}
          \resumeItemListStart
            \resumeItem{Geliştiricilerin canlı önizleme eşliğinde kolayca etkileşimli GitHub profil README dosyaları tasarlayıp temiz Markdown çıktısı alabileceği açık kaynaklı bir web uygulaması geliştirdim.}
          \resumeItemListEnd
      \resumeProjectHeading
          {\textbf{KıdıMusic Generator} $|$ \emph{TypeScript, JavaScript, CSS}}{\href{https://music.kidiman.com}{Canlı Demo} $|$ \href{https://github.com/yigitardakidiman/kidimusic-generator-app}{GitHub}}
          \resumeItemListStart
            \resumeItem{Metin yönlendirmeleriyle (prompt) özel müzik parçaları üreten, tarayıcıda anlık oynatma ve ses dosyası indirme imkanı sunan bir web uygulaması geliştirdim.}
          \resumeItemListEnd
      \resumeProjectHeading
          {\textbf{Offline Local RAG Asistanı} $|$ \emph{Python, RAG, Yerel Yapay Zeka}}{\href{https://github.com/yigitardakidiman/microsoft-rag}{GitHub}}
          \resumeItemListStart
            \resumeItem{Microsoft Foundry Local ve RAG kalıbı kullanarak, internet bağlantısı olmadan yerel dokümanlar üzerinden soru-cevap yapabilen çevrimdışı bir asistan geliştirdim.}
          \resumeItemListEnd
    \resumeSubHeadingListEnd

%-----------SERTİFİKALAR VE DİLLER-----------
\section{Sertifikalar \& Diller}
 \begin{itemize}[leftmargin=0.15in, label={}]
    \small{\item{
     \textbf{Sertifikalar}{: Python Programlama (BTK Akademi), HTML5 ile Web Geliştirme (BTK Akademi), Git \& GitHub ile Versiyon Kontrolü (BTK Akademi)} \\
     \textbf{Diller}{: Türkçe (Anadil), İngilizce (Orta Seviye - B1)}
    }}
 \end{itemize}

\end{document}
